\chapter{Application Architecture and the Module Framework}
\label{ch:arch}

\glance{OpenROAD is one C++ binary built from $\sim$35 modules under \file{src/}.
A single owner object, \id{ord::OpenRoad}, constructs every module once, injects
the shared services (the OpenDB database, the OpenSTA timer, and the
\id{utl::Logger}) into each, and exposes their commands through a Tcl/Python
(SWIG) layer. Read \file{src/OpenRoad.cc} once and you have the map of the whole
tool.}

\section{Entry point}

The binary starts in \file{src/Main.cc}: it parses command-line flags
(\cmd{-gui}, \cmd{-python}, \cmd{-exit}, \cmd{-metrics}, \cmd{-threads},
\cmd{-log}), creates the Tcl interpreter, constructs the application object, and
either runs an interactive shell or sources a script. The application object is
defined in \file{include/ord/OpenRoad.hh} and implemented in
\file{src/OpenRoad.cc}.

\id{ord::OpenRoad} is a process-wide singleton reached through the static
\id{OpenRoad::openRoad()}. It owns a pointer to every module and to the three
shared services, and hands them out through accessors such as \id{getDb()},
\id{getSta()}, and \id{getLogger()}.

\section{The three shared services}

Everything in OpenROAD is built on three objects that are created first and
passed to every module:

\begin{center}
\small
\begin{tabularx}{\linewidth}{@{}llX@{}}
\toprule
\textbf{Service} & \textbf{Type} & \textbf{Role} \\
\midrule
Database & \id{odb::dbDatabase} & the shared design + technology data (Ch.~\ref{ch:odb}) \\
Timer    & \id{sta::dbSta}      & static timing analysis (Ch.~\ref{ch:timing}) \\
Logger   & \id{utl::Logger}     & unified logging, metrics, debug (\pkg{src/utl}) \\
\bottomrule
\end{tabularx}
\end{center}

\section{Module construction and dependency injection}

\id{OpenRoad::init()} builds the whole tool. It creates the services, then
constructs each module, \emph{passing in the services and any sibling modules it
depends on}. This one function is the dependency graph of OpenROAD:

\begin{lstlisting}[style=orcode]
logger_    = new utl::Logger(log_filename, metrics_filename);
db_->setLogger(logger_);
sta_       = new sta::dbSta(tcl_interp, db_, logger_);
stt_builder_ = new stt::SteinerTreeBuilder(logger_);
opendp_    = new dpl::Opendp(db_, logger_);
global_router_ = new grt::GlobalRouter(logger_, service_registry_,
                                       stt_builder_, db_, sta_,
                                       antenna_checker_, opendp_);
resizer_   = new rsz::Resizer(logger_, db_, sta_, stt_builder_,
                              global_router_, opendp_, estimate_parasitics_);
tritonCts_ = new cts::TritonCTS(logger_, db_, getDbNetwork(), sta_,
                                stt_builder_, resizer_, estimate_parasitics_);
\end{lstlisting}

Note how, e.g., the resizer receives the timer, the Steiner builder, the global
router, and the detailed placer --- it \emph{uses} all of them. Table
\ref{tab:modclass} lists the module classes as constructed.

\begin{table}[h]
\centering
\small
\begin{tabularx}{\linewidth}{@{}llX@{}}
\toprule
\textbf{Class} & \textbf{Module} & \textbf{Stage} \\
\midrule
\id{ppl::IOPlacer}          & \pkg{src/ppl} & pin placement \\
\id{dpl::Opendp}            & \pkg{src/dpl} & detailed placement \\
\id{gpl::Replace}           & \pkg{src/gpl} & global placement \\
\id{mpl::MacroPlacer}       & \pkg{src/mpl} & macro placement \\
\id{par::PartitionMgr}      & \pkg{src/par} & partitioning \\
\id{cts::TritonCTS}         & \pkg{src/cts} & clock tree \\
\id{grt::GlobalRouter}      & \pkg{src/grt} & global route \\
\id{drt::TritonRoute}       & \pkg{src/drt} & detailed route \\
\id{ant::AntennaChecker}    & \pkg{src/ant} & antenna check \\
\id{rsz::Resizer}           & \pkg{src/rsz} & timing repair \\
\id{stt::SteinerTreeBuilder}& \pkg{src/stt} & Steiner trees \\
\id{est::EstimateParasitics}& \pkg{src/est} & RC estimation \\
\id{rcx::Ext}               & \pkg{src/rcx} & extraction \\
\id{pdn::PdnGen}            & \pkg{src/pdn} & power grid \\
\id{psm::PDNSim}            & \pkg{src/psm} & IR drop \\
\id{tap::Tapcell}           & \pkg{src/tap} & tap/endcap \\
\id{syn::Synthesis}         & \pkg{src/syn} & synthesis \\
\id{rmp::Restructure}       & \pkg{src/rmp} & ABC restructure \\
\id{cgt::ClockGating}       & \pkg{src/cgt} & clock gating \\
\id{dft::Dft}               & \pkg{src/dft} & scan/DFT \\
\id{fin::Finale}            & \pkg{src/fin} & metal fill \\
\id{dst::Distributed}       & \pkg{src/dst} & distributed compute \\
\id{ram::RamGen}            & \pkg{src/ram} & RAM generation \\
\bottomrule
\end{tabularx}
\caption{Module classes, as constructed in \id{OpenRoad::init()}.}
\label{tab:modclass}
\end{table}

\section{The per-module pattern}

Every module under \file{src/} has the same shape, so learning one teaches all:

\begin{itemize}
  \item \textbf{Public header} in \file{src/<mod>/include/<mod>/} defining the
        module's main class (e.g. \id{gpl::Replace}).
  \item \textbf{Implementation} in \file{src/<mod>/src/}.
  \item \textbf{A factory} \id{Make<Mod>.*} (e.g. \file{gpl/MakeReplace.h}) that
        \id{OpenRoad} calls to build the module and register its commands.
  \item \textbf{A SWIG interface} \file{<Mod>.i} and a \textbf{Tcl file}
        \file{<Mod>.tcl} that define the user commands and call into the C++ class.
  \item \textbf{A README.md} documenting every Tcl command and its options.
\end{itemize}

The template module \pkg{src/exa} (``example'') exists precisely as the
copy-me starting point for a new tool; \pkg{src/tst} holds shared test scaffolding.

\section{Commands: Tcl and Python via SWIG}

Tcl is the primary interface; Python is secondary. The top-level wrappers are
\file{src/OpenRoad.i}, \file{src/OpenRoad.tcl}, and the Python variants
\file{src/OpenRoad-py.i} / \file{src/Design.i}. Each module's \file{*.i} is
compiled by SWIG into Tcl and Python bindings; the hand-written \file{*.tcl}
adds argument parsing and calls the SWIG-wrapped C++ methods.

The Python-facing convenience classes live in \file{src/Design.cc},
\file{src/Tech.cc}, and \file{src/Timing.cc} (the \id{Design}, \id{Tech}, and
\id{Timing} objects used in \cmd{openroad -python}).

\paragraph{Finding a command's code.} Grep the command name in the module's
\file{*.tcl} (Tcl argument handling) and \file{*.i} (SWIG binding); both point
into the module's C++ class. The \id{ord::} namespace also provides helpers such
as \id{ord::get\_db} and \id{ord::get\_db\_block}.

\section{Logging and metrics: \pkg{src/utl}}

All output goes through \id{utl::Logger} (never raw \id{printf}/\id{cout}). Each
message carries a \emph{tool ID} and a numeric code, giving stable identifiers
like \texttt{CTS-0001} that the \cmd{man} command can explain.

\begin{center}
\small
\begin{tabularx}{\linewidth}{@{}lX@{}}
\toprule
\textbf{Call} & \textbf{Use} \\
\midrule
\id{report()} & always-printed user output \\
\id{info(tool, id, ...)} & informational message with stable ID \\
\id{warn(tool, id, ...)} & warning \\
\id{error(tool, id, ...)} & recoverable error (throws) \\
\id{critical(tool, id, ...)} & fatal error (aborts) \\
\id{metric(name, value)} & emit a JSON metric (\cmd{-metrics}) \\
\id{debug(tool, group, ...)} & gated debug stream \\
\bottomrule
\end{tabularx}
\end{center}

\pkg{src/utl} also provides the \cmd{man} help command, the \cmd{tee} output
redirect, and an optional Prometheus metrics endpoint
(\id{utl::startPrometheusEndpoint}).

\glance{To add a tool: copy \pkg{src/exa}, define your class + \id{Make} factory,
add a \file{.i}/\file{.tcl} command pair, construct it in \id{OpenRoad::init()},
and log through \id{utl::Logger}. That is the entire contract.}
